\section{Previous work}

\subsection{Embeddable Common Lisp (ECL)}

ECL is implemented as a core written in \clanguage{}.  It is built by
first compiling this core and then adding code written in \cl{} by
loading it into the system resulting from the compilation of the core.
For that reason, the reader must be written largely in \clanguage{}. 

When an attempt is made to read an incomplete expression, so that the
\texttt{end-of-file} condition is invoked, no restart other than
\texttt{restart-toplevel} is available.

Similarly when an attempt is made to read a symbol with an explicit
package marker where the corresponding package does not exist, an
error of type \texttt{simple-error} is signaled, and again, the only
restart available is \texttt{restart-toplevel}.

When an attempt is made to read an incorrect token in the form of a
malformed symbol with three package markers, again an error of type
\texttt{simple-error} is signaled, and the message is cryptic.  For
example, an attempt to read the malformed symbol \texttt{cl:::car},
the error message is ``There is no package with the name CL:::C''.
